\section*{Theorem 7 --- the $D_5\times A_3$ lift of the $E_8$ atlas}
% =====================================================================
\addcontentsline{toc}{section}{Theorem 7}
Since $A_3=\mathfrak{su}_4\cong\mathfrak{so}_6=D_3$, one has the standard chain
$E_8\supset \mathfrak{so}_{16}\supset \mathfrak{so}_{10}\oplus\mathfrak{so}_6=D_5\times A_3$, with the
exact branching (verified dimensions and root/weight counts):
\begin{equation}
\begin{aligned}
  &\boxed{\ \mathfrak e_8=(45,1)\oplus(1,15)\oplus(10,6)\oplus(16,4)\oplus(\overline{16},\overline4)\ },\\
  &\underbrace{45{+}15{+}60{+}64{+}64}_{=248}=\dim E_8,\qquad
  \underbrace{40{+}12{+}60{+}64{+}64}_{=240}=|R(E_8)| .
\end{aligned}
\end{equation}
The adjoint $\mathbf{120}$ of $\mathfrak{so}_{16}$ is $(45,1)\oplus(1,15)\oplus(10,6)$ and its spinor
$\mathbf{128}$ is $(16,4)\oplus(\overline{16},\overline4)$. The TFPT dictionary makes each summand a
known object:
\begin{center}\small
\begin{tabularx}{\linewidth}{@{}llY@{}}
\toprule
piece & value & TFPT reading \\
\midrule
$(45,1)$ & $45=40+5$ & $D_5$ carrier adjoint: Clebsch edges $40=|R(D_5)|$ plus carrier rank $\gcar=5$ \\
$(1,15)$ & $15=12+3$ & $A_3$ family$+$Higgs adjoint: $\dimg=12$ plus family rank $\Nfam=3$ \\
$(10,6)$ & $10\cdot6$ & $\Lambda$ decuple $\AL=\binom52=10$ times $Z_6$ transport $\binom42=6$ \\
$(16,4)\oplus(\overline{16},\overline4)$ & $2\cdot64$ & matter half-spinor $S^+$ times the four corners $\mu_4$, plus conjugate sheet \\
\bottomrule
\end{tabularx}
\end{center}

\noindent\textbf{Proof route (replaces the old blind $240$-metrization).}
\emph{Stage 1 ($D_5$, \mE):} $C^+$ is the $D_5$ half-spin orbit, $R_4=$ Clebsch, $\operatorname{Aut}=W(D_5)$,
$|E(R_4)|=40=|R(D_5)|$ --- Theorem 2.
\emph{Stage 2 ($A_3$, \mE):} the carrier geometry fixes $X_f^\circ=\mathbb P^1\setminus\mu_4$ with
$H_1(X_f^\circ,\mathbb Z)=\mathbb Z^3$ and a rigid $D_4\subset PGL_2$ generated by $z\mapsto iz$ and
$z\mapsto 1/z$. The logarithmic forms $\omega_i=dz/(z-p_i)$ regular at infinity span $H^1$ by their
differences, dual under the residue pairing to the puncture loops $\gamma_i$ with $\operatorname{Res}(\omega_i,\gamma_j)=\delta_{ij}$;
hence the canonical (residue) metric on the puncture classes is $I_4$, and the twelve
puncture-difference cycles $\gamma_i-\gamma_j$ are exactly the $A_3$ roots:
\begin{equation}
  \boxed{\ \begin{aligned}
    &\{\gamma_i-\gamma_j\}_{i\ne j}:\ \#=12=|R(A_3)|=\dimg,\quad \operatorname{rank}=3=\Nfam,\\
    &\text{Gram}=\text{Cartan}(A_3),\ \det=4
  \end{aligned}\ }\qquad\mE ,
\end{equation}
with $\dim A_3=15=\dim S^+-1$, and the $A_3$ positive roots equal the exponent sum $1{+}2{+}3=6=Z_6=|E(K_4)|$.
\emph{The corner rotation realises the $A_3$ Coxeter element.} The carrier generator $z\mapsto iz$
permutes $\mu_4$ as the $4$-cycle $(p_1p_2p_3p_4)\in S_4=W(A_3)$ and acts on $H^1=\mathbb C^3$ as an
isometry of the Cartan form with
\begin{equation}
  \boxed{\ \text{ord}=4=h(A_3)=|\mu_4|,\quad \text{spec}=\{i,-1,-i\},\quad
  \chi(t)=\tfrac{t^4-1}{t-1}=\Phi_4(t)\Phi_2(t)\ }\qquad\mE ,
\end{equation}
i.e.\ the eigenvalues are the $A_3$ exponents $\{1,2,3\}$ as fourth roots of unity --- the exact analogue
of the $E_8$ Coxeter element reading $\Phi_{30}$. The full $D_4=\langle z\mapsto iz,z\mapsto1/z\rangle$
embeds in $W(A_3)=S_4$ (order $8$, index $3$) by isometries. \emph{Stage 3 (standard, \mE):} assemble
$E_8$ from the $D_5\times A_3$ branching --- no $240$-map required.

\noindent\textbf{The $\mu_4$ glue theorem (why the four corners appear everywhere).} The two factors
have equal discriminant groups, and $E_8$ is unimodular:
\begin{equation}
  \boxed{\ \operatorname{disc}(D_5)=\operatorname{disc}(A_3)=\mathbb Z_4,\qquad
  [\,E_8:D_5\oplus A_3\,]=\sqrt{\tfrac{\det D_5\cdot\det A_3}{\det E_8}}=\sqrt{\tfrac{4\cdot4}{1}}=4=|\mu_4|\ }\qquad\mE .
\end{equation}
So $E_8=(D_5\oplus A_3)+\mu_4$-\emph{glue}: the four family corners $\mu_4=\mathbb Z_4$ are precisely the
glue code that fuses the carrier and the family into the unimodular $E_8$ lattice. The glue pairs the
$D_5$ spinor class ($=S^+$, the family) with the $A_3$ class --- exactly the $(16,4)\oplus(\overline{16},\overline4)$
spinor block. Hence $\mu_4$ plays three roles at once: family corners, $A_3$ origin, and $E_8$ glue.

\noindent\textbf{Discriminant-norm glue lemma (and the TFPT $5/4,3/4$ identification).} The two glue
generators carry the discriminant quadratic form $q\in\mathbb Q/2\mathbb Z$:
\begin{equation}
  \boxed{\ \begin{aligned}
    &q(D_5)=\tfrac{n}{4}\big|_{n=5}=\tfrac54\ (\text{spinor class}),\quad
    q(A_3)=\tfrac{n}{n+1}\big|_{n=3}=\tfrac34\ (\text{corner class}),\\
    &q(D_5)+q(A_3)=2\equiv0\!\!\pmod 2,\qquad q(D_5)-q(A_3)=\tfrac12
  \end{aligned}\ }\qquad\mE .
\end{equation}
The diagonal $\langle(s,g)\rangle\cong\mathbb Z_4$ is isotropic ($k^2\cdot2\equiv0$ for all $k$), so the
glue closes to an \emph{even unimodular} lattice --- this is $E_8$.

\textbf{The two glue norms are just $(\gcar,\Nfam)$ over the glue index.} Because $\operatorname{rank}D_5=\gcar=5$
(so $q(D_5)=\tfrac n4|_{n=5}=\gcar/|\mu_4|$) and $\operatorname{rank}A_3=\Nfam=3$
(so $q(A_3)=\tfrac{n}{n+1}|_{n=3}=\Nfam/|\mu_4|$, with $n{+}1=4=|\mu_4|$),
\[
  \boxed{\ q(D_5)=\frac{\gcar}{|\mu_4|}=\tfrac54,\qquad q(A_3)=\frac{\Nfam}{|\mu_4|}=\tfrac34\ }
\]
and \emph{all four arithmetic operations} on the pair are then forced compiler atoms:
\[
  \begin{aligned}
  &\underbrace{q(D_5){+}q(A_3)=\tfrac{\gcar+\Nfam}{|\mu_4|}=2}_{\text{$E_8$ root norm}}\,,\qquad
  \underbrace{q(D_5){-}q(A_3)=\tfrac{|\Z_2|}{|\mu_4|}=\tfrac12}_{\delta\ \text{(lepton transport)}}\,,\\[4pt]
  &\underbrace{q(D_5)q(A_3)=\tfrac{\gcar\Nfam}{|\mu_4|^2}=\tfrac{15}{16}}_{\dim\mathfrak{su}(4)/\dim S^+}\,,\qquad
  \underbrace{\tfrac{q(D_5)}{q(A_3)}=\tfrac{\gcar}{\Nfam}=\tfrac53}_{}\,.
  \end{aligned}
\]
So the distinguished lepton transport value $\delta=\tfrac12=(\gcar-\Nfam)/|\mu_4|=|\Z_2|/|\mu_4|$ is the
carrier-minus-family count over the glue index --- a fully atomic origin, not a chosen rational (it also
equals the harmonic-frame holonomy diagonal modulus, \veri{v41\_leg\_assignment.py}); and $(\text{sum},\text{ratio})$
invert uniquely back to $(\tfrac54,\tfrac34)$. \veri{v51\_boundary\_half\_step.py} The sum $\tfrac54+\tfrac34=2$ is
literally the $E_8$ spinor-root norm $\tfrac12(\pm)^8$ split across $D_5\times D_3$:
$5\cdot\tfrac14+3\cdot\tfrac14=2$. Strikingly, both norms are pre-existing TFPT constants:
\begin{equation}
  \boxed{\ \frac{\dtwo}{\dtop^2}=\frac{2880\,\cthree^8}{(48\,\cthree^4)^2}=\frac54=q(D_5),\qquad
  \xi_{\mathrm{tree}}=\frac{\dimg}{\dim S^+}=\frac{12}{16}=\frac34=q(A_3)\ }\qquad\mE ,
\end{equation}
so $\dtwo/\dtop^2+\xi_{\mathrm{tree}}=2$: the two-loop/defect ratio is the $D_5$ spinor norm and the
gravitational tree normalization is the $A_3$ corner norm, and their sum is the $E_8$ root norm.

\begin{keybox}[The whole integer skeleton from one pair $(\gcar,\Nfam)=(5,3)$ \mE\ (\veri{v53\_compiler\_core.py})]
Every load-bearing integer of the compiler flows from the single boundary pair
$(\gcar,\Nfam)=(5,3)$ by sum, difference and squares:
\[
  \operatorname{rank}E_8=\gcar{+}\Nfam=8,\qquad
  |\Z_2|=\gcar{-}\Nfam=2,\qquad
  |\mu_4|=\tfrac{\gcar+\Nfam}{2}=4,
\]
and $(\Nfam,|\mu_4|,\gcar)=(3,4,5)$ is \emph{the} Pythagorean triple. This is not decorative: the
grand-mass-volume exponent $\Delta_Y=\gcar^2=25$ ($\det M_{\mathrm{SM}}\sim(\phiz)^{25}$; document~3)
splits over the sectors as a \emph{difference of squares},
\[
  \boxed{\ \Delta_Y=\gcar^2
  =\underbrace{\Nfam^2}_{\text{down}=9}
  +\underbrace{(\gcar{-}\Nfam)(\gcar{+}\Nfam)}_{=\,|\Z_2|\cdot\operatorname{rank}E_8}
  =9+16=\dim S^+\ }
\]
i.e.\ the down sector carries $\Nfam^2$ and up$+$lepton carries $|\Z_2|\!\cdot\!\operatorname{rank}E_8=\dim S^+$
(the $K$-row sums $(6,9,10)$). \textbf{The anchor is the Dirac square root of this skeleton:} $a=(1,1,2)$ is
the \emph{unique} $3$-multiset whose elementary symmetric functions are $(|\mu_4|,\gcar,|\Z_2|)=(4,5,2)$,
equivalently its characteristic polynomial is
$\chi_a(t)=(t{-}1)^2(t{-}2)=t^3-|\mu_4|\,t^2+\gcar\,t-|\Z_2|$ --- the compiler atoms \emph{are} its
char-poly coefficients. Sharper still (\veri{v491\_p2\_partition\_corollary.py}): the partition needs only
$e_1=4$ --- three \emph{positive integers} summing to $4$ already force $\{1,1,2\}$, so $\gcar=e_2=5$ and
$|\Z_2|=e_3=2$ are \emph{corollaries} of the four marks given the weight typing (integrality $=$
Birkhoff--Grothendieck splitting exponents, positivity $=h^0{=}0$, sum $=$ par-deg-$0$ ``given (U)'');
every dropped assumption is negative-controlled, and \texttt{AX.P2.01} stays a declared axiom --- the
residual is the typing, not the integer $5$. That typing is itself now hardened
(\veri{v499\_p2\_weights\_deligne\_bg.py}, \texttt{P2.TYPING.01}): given the four marks, rank $3$, the cusp
class and (U), Deligne $+$ Birkhoff--Grothendieck $+$ Mehta--Seshadri stability force
$E=\mathcal O(-2)\oplus\mathcal O(-1)^2$ outright ($h^0{=}0$ is the positivity), leaving only the module
identification and (U) as \mC\ residuals --- and the module-identification rest is now \emph{one}
residual with the QGEO modulus rest, both reducing to the same order-$4$ clock $=$ Coxeter/cusp-class
carrier datum (\vref{v503}, \texttt{QGEO.EMERGE.LIGHT.01}; markers unchanged); the clock-rigidity
theorems then reduce that common residual to one alignment bit, with the clock's order fermionically
forced, $U^2=(-1)^F$ exactly (\veri{v506\_seam\_clock\_rigidity.py}, \texttt{SEAM.CLOCK.RIGIDITY.01});
and the bit-origin theorems derive the deck's \emph{class} (a half-period translation of the seam
double) --- only its position stays the \mC\ carrier input
(\veri{v507\_seam\_bit\_origin.py}, \texttt{SEAM.BIT.ORIGIN.01}); the freedom theorems then type
that position half as \emph{topology} (the covering deck is free on the seam circle, the edge class
is excluded), so the residual is the square-modulus datum $\tau=i$ alone
(\veri{v510\_seam\_bit\_freedom.py}, \texttt{SEAM.BIT.FREEDOM.01}); and the flag-transitivity
equivalence web restates that residual in discrete form --- flag transitivity of the four marks
($V_4\to D_4$) $\iff\tau=i$, with bare mark-transitivity automatic for every configuration and no
local jet seeing the side bit (\veri{v512\_seam\_tau\_flag.py}, \texttt{SEAM.TAU.FLAG.01}) ---
a residual that has since survived ten side-blind derivation attacks
(\vref{v521}/\vref{v525}/\vref{v529}) and carries a physical definition as the twist-class
choice with a gauge-robust order parameter, in principle interferometrically readable
(\vref{v528}; it remains formal input) --- and, since \vref{v534}, the first \emph{dynamical}
selection datum: reflection positivity on the interacting mark family keeps exactly the
$\delta=\pi/2$ member alive (toy-fenced).
So the theory's only transcendental is $\pi$ (through
$\cthree=\tfrac1{8\pi}$, $8=\operatorname{rank}E_8$, and the seed $\tfrac1{6\pi}$, $6=|R^+(A_3)|$);
everything else is this $(5,3)$-generated integer skeleton.
\end{keybox}

\begin{keybox}[Boundary Carrier Selection Theorem (Theorem A) \mE/\mC]
$E_8$ is not assumed: it is \emph{selected} as the unique even-unimodular audit hull of the admissible
boundary-carrier data. \textbf{Boundary side:} $\mu_4=\{1,i,-1,-i\}$ gives $X_f=\PP^1\setminus\mu_4$ with
$H_1=\Z^3$, hence $A_3$ and $\Nfam=3$. \textbf{Carrier side:} a minimal $D$-type carrier with a $16$-dim
half-spinor forces $D_5$ (only $n=5$ gives $2^{n-1}=16$; $S^+=\Lambda^{\mathrm{even}}\mathbb C^5=1{+}10{+}5$).
\textbf{Glue:} $\operatorname{disc}(D_5)=\operatorname{disc}(A_3)=\Z_4$ (isomorphic, cyclic, anti-isometric),
$q(D_5)+q(A_3)=2$ (the $E_8$ root norm), glue index $4=|\mu_4|$, and $E_8$ is the unique even unimodular
rank-$8$ lattice. \textbf{Exclusion:} the four conditions
\[
  \text{(i) }\dim S^+{=}16,\quad\text{(ii) }\Nfam{=}3,\quad\text{(iii) cyclic }\Z_4\text{ glue},\quad
  \text{(iv) }q_D{+}q_A{=}2
\]
are met \emph{only} by $D_5\oplus A_3$; $D_8$ (non-cyclic $\Z_2{\times}\Z_2$, no family geometry),
$E_7{\oplus}A_1$ ($1$ family, no $16$-spinor), $E_6{\oplus}A_2$ ($2$ families, $\Z_3$), $A_4{\oplus}A_4$
(no $D$-spinor, $\Z_5$) and $A_8$ (single factor) each fail at least one. \textbf{Typing:} the
lattice/classification facts are \mE; the physical selection (reflection-positive seam $\cthree=1/8\pi$ +
minimal chiral carrier $\Rightarrow$ exactly these inputs) is the \mC/\mO\ selection axiom (P1 analytic
interface $+$ P2 carrier interface), not a zero-input derivation. \veri{v47\_selection\_theorem.py}
\end{keybox}

\noindent\textbf{Group-level closure ($\mu_4$ as the common centre).} The lattice glue lifts to a clean
group statement, because both factors share the \emph{same} centre and $\mu_4$ \emph{is} the diagonal:
\begin{equation}
  \boxed{\ Z(\mathrm{Spin}(10))=\mathbb Z_4,\quad Z(SU(4))=\mathbb Z_4,\quad |\mu_4|=4,\quad
  G_{\mathrm{closure}}=\frac{\mathrm{Spin}(10)\times SU(4)}{\Delta\mathbb Z_4}\subset E_8\ }\qquad\mE ,
\end{equation}
with $\operatorname{rank}=5+3=8=\operatorname{rank}E_8$. Every adjoint block is diagonal-$\mathbb Z_4$
neutral (generator $\omega=i$; charges mod $4$): $(45,1),(1,15)\!:\!0$; $(10,6)\!:\!2{+}2{=}0$;
$(16,4)\!:\!1{+}3{=}0$; $(\overline{16},\overline4)\!:\!3{+}1{=}0$, so $45{+}15{+}60{+}64{+}64=248$ are
genuine reps of the quotient. This promotes $\mu_4$ from a lattice glue index to the \emph{common centre
quotient} --- the physically cleanest language for the gluing: the $D_5$ half-spinor $16$ pairs with the
$SU(4)$ fundamental $4$ so the spinor charges $\pm1$ cancel the fundamental charges $\mp1$.

\noindent\textbf{Root-length completion.} Every $E_8$ block reaches norm $2$:
\begin{center}\small
\begin{tabularx}{\linewidth}{@{}llY@{}}
\toprule
block & norm split & reading \\
\midrule
$(45,1),(1,15)$ & $2,\ 2$ & $D_5$, $A_3$ root adjoints \\
$(10,6)$ & $1+1=2$ & $SO(10)$ vector $\oplus$ $SO(6)$ vector \\
$(16,4)\oplus(\overline{16},\overline4)$ & $\tfrac54+\tfrac34=2$ & $D_5$ spinor norm $+$ $A_3$ corner norm \\
\bottomrule
\end{tabularx}
\end{center}
The old question ``can the $240$ flip atlas be metrized to $E_8$?'' is now the sharp question ``are
the TFPT $D_5$ and $A_3$ metrics in standard normalization?'' --- and the $5/4,3/4$ match says yes.

\noindent\textbf{Explicit construction (the lift is now built, not counted).} Embedding
$D_5\oplus A_3=D_5\oplus D_3\subset D_8$ and adjoining the even spinor coset (the $\mu_4$ glue), the
$240$ roots are produced explicitly and verified:
\begin{equation}
  \boxed{\ \{\pm e_i\pm e_j\}_{i<j}^{8}\ (112)\ \cup\ \{\tfrac12(\pm)^8:\ \text{even}\}\ (128)
  \ =\ 240\ \text{roots, all norm }2,\ \text{Cartan }\det=1\ }\qquad\mE .
\end{equation}
Splitting each root into its $D_5$ (first five) and $A_3$ (last three) parts gives exactly the
branching blocks --- $(2,0)\!:\!40$, $(0,2)\!:\!12$, $(1,1)\!:\!60$, and $(\tfrac54,\tfrac34)\!:\!128$ ---
and the $128$ spinor roots realise $\tfrac54+\tfrac34=2$ on the nose. All $240$ vectors are integer
combinations of one set of $8$ simple roots, so the assembled lattice \emph{is} $E_8$. The blind
$240$-metrization is therefore eliminated: $E_8$ is explicitly $(D_5\oplus A_3)+\mu_4$-glue.
\begin{remark}[the metric, computed]
This reframes $E_8$ from origin to closure: $E_8$ is the closure algebra of the $D_5$ carrier and the
$A_3$ family geometry, glued by $\mu_4$. The root system, Cartan form, $W(A_3)=S_4$ and the
corner-rotation $=$ $A_3$ Coxeter element are exact (residue pairing $+$ explicit carrier symmetry).
The transcendental geometric metric $G_{ij}=-\log|p_i-p_j|$ on $\mu_4$ has been computed: it is only
$D_4$-symmetric, with a regularisation-independent anisotropy $-\log2$ on the root space, hence
\emph{not} proportional to $\mathrm{Cartan}(A_3)$. It need not be: the $E_8$-relevant $A_3$ is the
canonical \emph{residue} pairing, and the geometric metric is a $D_4$-equivariant deformation that
preserves the root system, Weyl group and Coxeter element. So the lattice data is exact \mE; the
analytic metric is a \emph{characterised} deformation, not an open gap.
\end{remark}

\noindent\textbf{Uniform Coxeter law.} All three factors obey $|R|=h\cdot\operatorname{rank}$, and the
two TFPT Coxeter numbers are themselves carrier data:
\begin{equation}
  \boxed{\ \begin{aligned}
    &h(D_5)=8=r(E_8)=\varphi(30),\quad h(A_3)=4=|\mu_4|,\\
    &h(D_5)\,h(A_3)=2^{\gcar}=32,\quad 2^{\gcar}-2=30=h(E_8)
  \end{aligned}\ }\qquad\mE ,
\end{equation}
so the Coxeter closure of Theorem 3 reappears as $h(E_8)=h(D_5)h(A_3)-2$, and
$h(D_5)+h(A_3)=12=\dimg=|R(A_3)|$.

% =====================================================================
\section*{Theorem 8 --- the $(1,1,2)$ anchor microcode and the $\varphi_0$-ladder}
% =====================================================================
\addcontentsline{toc}{section}{Theorem 8 --- anchor microcode}
\providecommand{\phiz}{\varphi_0^{\mathrm{ret}}}
The parabolic flavor bundle $E_\bullet=\mathcal O(-2)\oplus\mathcal O(-1)^2$ on $\mathbb P^1\setminus\mu_4$
has splitting type $a=(1,1,2)$, whose power sums $p_n(a)=2+2^n$ generate the whole compiler budget:

\par\medskip\noindent\textbf{The anchor power compiler.}
\[
  (p_1,\dots,p_5)=(4,6,10,18,34)=\bigl(|\mu_4|,\,|R^+(A_3)|,\,\AL,\,\Nfam|R^+(A_3)|,\,\text{$\mathbb Z_6$ lift}\bigr),
\]
\[
  q(A_3)=\tfrac{p_2}{2p_1}=\tfrac34,\quad q(D_5)=\tfrac{p_3}{2p_1}=\tfrac54,\quad \textstyle\sum L=p_1p_3=40,
\]
\[
  \Oadm=2p_1p_2=48,\qquad D_{\mathrm{start}}=p_2p_3=60,\qquad |R(E_8)|=4p_2p_3=240.
\]
The elementary symmetric data $\chi_a(t)=(t-1)^2(t-2)$ give $(e_1,e_2,e_3)=(4,5,2)=(|\mu_4|,\gcar,|\mathbb Z_2|)$,
so the anchor reconstructs the carrier rank, and $10b_1=e_1^2+e_2^2=4^2+5^2=41$. The seed itself is the
two-term anchor form $\phiz=\tfrac{2p_1}{p_2}\cthree+2p_1p_2\,\cthree^4=\tfrac43\cthree+48\cthree^4$.

\begin{keybox}[The anchor affine normal form (2026-08-07) --- all budgets from ONE recursion \mE\ (\veri{v832\_anchor\_flavor\_checksum.py})]
The anchor power compiler above is one affine map.  The power sums $p_n(a)=2+2^n$ satisfy
\[
  \boxed{\ p_{n+1}=T(p_n),\qquad T(x)=2x-2\ }\qquad\text{identically in $n$ (sympy),}
\]
with the shift $q_n=p_n-2=2^n$ pure binary doubling.  $T$ has the \emph{unique} fixed point
$2=|\Z_2|$ (the sheet), and the $T$-orbit of $4$ is exactly the compiler quintet
$(4,6,10,18,34)=(|\mu_4|,\,|R^+(A_3)|,\,\AL,\,\Nfam|R^+(A_3)|,\,\text{$\Z_6$ lift})$, with
$p_0=3=\Nfam$ prepended by $T(3)=4$.  Every budget line is then a word in the orbit and the fixed
point: $|R(E_8)|=p_1p_2p_3=240$; $\dim E_8=240+q_3=248$, where the ladder identity
$p_4-p_3=p_3-2=8=h(D_5)=\operatorname{rank}E_8$ makes the corpus form $240+(p_4{-}p_3)$ the
\emph{same} formula; $h(E_8)=p_2p_3/2=30$ (so $D_{\mathrm{start}}=p_2p_3=60$); $|R(D_5)|=p_1p_3=40$;
$\Oadm=2p_1p_2=48$; and the elementary layer $(e_1,e_2,e_3)=(4,5,2)$ has the fixed point of $T$ as
$e_3$, with $10b_1=e_1^2+e_2^2=41$.  Controls: the wrong map $T'(x)=2x-1$ (fixed point $1$, orbit
misses the quintet) and the wrong anchor $a'=(1,2,2)$ (recursion offset $-1$, budget misses
$(240,48)$) both fire.  A normal form of the existing compiler tables --- a compression, not a new
claim; $18/18$ exact, verdict \textsc{anchor-affine-exact}.
\end{keybox}

\noindent\textbf{The $\varphi_0$-ladder.} With $\lambda_Y=\sqrt{\phiz(1-\phiz)}$, every fermion mass is the
single closed form $\hat m_{f,j}=\tfrac{v_{\mathrm{geo}}}{\sqrt2}\lambda_Y^{L_{f,j}}\Lambda_{f,j}$ with
$\lambda_Y^{L}\Lambda=\pi\,c\,(\phiz)^{k}$, so all nine holonomies are resolvent outputs and the masses
form the $\phiz$-order ladder $t(0)>\{b,c,\tau\}(2)>\{s,\mu\}(3)>\{u,d\}(4)>e(5)$.

\noindent\textbf{Secondary $E_8$ atlas.} Beyond the main $D_5\times A_3$ split, the residue matrix appears
in the $E_6\times A_2$ branching, $248=\|R\|_F^2+\det R+2(\mathbf 1^\top Ra)\Nfam=78+8+2\cdot27\cdot3$, and
the scalaron exponent in $E_7\times A_1$, $112=7\cdot16$. The seven $E_8$ slices are charted in the
companion note \texttt{e8\_secondary\_branching\_atlas} as a falsifiable audit raster.
\par\smallskip\noindent\textbf{The $E_6{\times}A_2$ block is the $\tau{=}\omega$ face
(\veri{v407\_dn\_pairings\_omega.py}, \texttt{SEAM.EQUIV.02}).} $A_2$ is the Eisenstein lattice of the family
CM point $\tau{=}\omega$ (Cartan $\det=3=\Nfam$), and the $(d,n)$ selector's three pairings
$\{n{\cdot}\mathbf1,\,n{\cdot}a,\,n{\cdot}\sigma\}=\{2,8,121\}$ are exactly its family-slice atoms
$\{|Z_2|,\ \dim A_2{=}\det R{=}\operatorname{rank}E_8,\ \|Pl(K)\|_1^2\}$, so deriving $(d,n)$ (the R4$'$
residue, \vref{v139}) folds into the dual keystone (\vref{v405}), leaving only ``why $11{=}\|Pl(K)\|_1$''.

% =====================================================================
\section*{Master statement and status}
% =====================================================================
\addcontentsline{toc}{section}{Master statement}
\begin{keyeq}
Five carrier slots give $D_5$,\ four family corners give $A_3$,\ and $D_5\times A_3$ gives the $E_8$ atlas.
\end{keyeq}
\noindent Equivalently: $\AL=\tfrac{h}{\Nfam}=10=\binom52$, $h=30=2\cdot3\cdot5$, $(5,3,2):\ \sum=10,\ \prod=30$;
$\Lambda$, $\det M_\ell$, the $E_8$ Coxeter structure and the action ladder read the same five-slot
pair-decuple. The Standard model, flavor, $\Lambda$, inflation and $E_8$ are projections of the
even carrier code $C^+$ ($=D_5$ half-spinor) on five slots together with the four-point family $A_3$.

\begin{center}\small
\begin{tabularx}{\linewidth}{@{}p{0.52\linewidth}p{0.10\linewidth}Y@{}}
\toprule
Statement & Status & Note \\
\midrule
Even code $C^+$, $\dim S^+=16$, dual-root $Y$ & \mE & Theorem 1 \\
Exterior--Action duality $(1,10,5)\!\leftrightarrow\!(1,5,10)$ & \mE & scheme eigenvalues verified \\
$R_4=$ Clebsch $(16,5,0,2)$, $R_2=(16,10,6,6)$; $\operatorname{Aut}_{\mathrm{aff}}=W(D_5)=1920$ & \mE & carrier $=D_5$ half-spinor \\
$2^g-2=g(g+1)\Leftrightarrow g=5$ & \mE & Coxeter closure (unique) \\
$h{=}30{=}2{\cdot}3{\cdot}5$, $r{=}\varphi(30){=}8$, $\chi_C{=}\Phi_{30}$ & \mE/\mC & cyclotomic compiler \\
$|R^+|=\binom{16}{2}=\operatorname{Tr}X^2=4h=120$ & \mE & positive-edge moment \\
$40,60,120,240=10\cdot(4,6,12,24)$; $41{=}1{+}4{\cdot}10$ & \mE & $K_5\times K_4$ product table \\
$10b_1{=}5r{+}1$, $u$ totient, $M_{\mathrm{scal}}^2{=}\cthree^{r-1}$, $\omega_b{=}\lC/10$ & \mE/\mC & downstream readouts \\
$A_s=N_\star^2\cthree^7/24\pi^2$ & \mC/\mO & if $M_{\mathrm{scal}}=M_{\mathrm{Star}}$ \\
$\mathfrak e_8=(45,1)\oplus(1,15)\oplus(10,6)\oplus(16,4)\oplus(\overline{16},\overline4)$ & \mE/\mC & $D_5\times A_3$ lift \\
$240$ roots built from $D_5{\oplus}A_3{+}\mu_4$ glue, $\det 1$ & \mE & explicit $E_8$ construction \\
flavor $\sum L{=}40{=}|R(D_5)|$, entries on $C_6{=}|R^+(A_3)|$ & \mE & transport $=D_5{\times}A_3$ \\
$A_3$ from $\mathbb P^1\setminus\mu_4$: residue pairing $=$ Cartan, $\det 4$ & \mE & Stage 2 (Theorem 7) \\
corner $z\mapsto iz=$ $A_3$ Coxeter el.: $\{i,-1,-i\}$, $\Phi_4\Phi_2$ & \mE & order $4=h(A_3)=|\mu_4|$ \\
$\operatorname{disc}(D_5)=\operatorname{disc}(A_3)=\mathbb Z_4$, glue $=\mu_4$ & \mE & $E_8=(D_5{\oplus}A_3){+}\mu_4$ \\
$q(D_5){=}\tfrac54{=}\dtwo/\dtop^2$, $q(A_3){=}\tfrac34{=}\xi_{\mathrm{tree}}$, sum $2$ & \mE & discriminant-norm glue \\
$120{=}\dim\mathfrak{so}_{16}$; $\Oadm{=}40{+}1{+}7$ & \mE & adjoint half; deficit ladder \\
$h(D_5){=}8{=}r(E_8)$, $h(A_3){=}4{=}|\mu_4|$, $|R|{=}h\cdot r$ & \mE & uniform Coxeter law \\
$\operatorname{Exp}(E_8)$ as $41/48$ budget deficits & \mO & exponent ledger (fingerprint) \\
geometric metric $D_4$-deform (anisotropy $-\log2$), residue $=A_3$ & \mE & metric computed/characterised \\
\bottomrule
\end{tabularx}
\end{center}

\noindent\emph{Scope.} This note compresses; it does not re-derive the upstream boundary kernel
(the P1 hardening chain) or the carrier rigidity theorem (Lean), and it makes no claim of a final theory. The
arithmetic and the $D_5$/Clebsch/$W(D_5)$, $A_3$/residue-Cartan/Coxeter, and $\mu_4$-glue
identifications are exact, and the $E_8$ lattice is now built explicitly from $D_5\oplus A_3$ plus the
$\mu_4$ glue ($240$ roots, $\det1$). The only residual \mC\ is the analytic identification of the carrier geometry's
$L^2$ cusp Hodge metric with the canonical residue metric up to scale, which does not affect the
explicit root/lattice/Weyl construction. Honest non-fits (Koide $0.6645\neq2/3$, $\eta_B$, $m_p/m_e$)
stay outside the compiler.

% =====================================================================
